\section{Meganekko: интеллект, наблюдение, сдержанность, nerd-код}

Meganekko — интеллект как визуальный код: очки здесь не аксессуар, а сигнал наблюдательности и сдержанности, тот же принцип усиленного символа из 0.2, применённый к одной детали лица. Нодока Миядзаки из "K-On!" — скромный, но точный пример: единственная в компании персонажей-музыкантов, кто держит дистанцию наблюдателя, а не участника хаоса, и очки визуально маркируют эту позицию раньше, чем она успевает что-то сказать.

Функция meganekko в группе обычно структурная, а не самостоятельная: она точка, с которой сцена смотрит на остальных персонажей со стороны, комментирует, замечает то, что участники происходящего не видят изнутри. Это близко к позиции хора из 1.3 — голос, который проговаривает то, что зритель и так подозревал.

Баг этого амплуа тоже структурный, а не случайный: наблюдатель, который хорошо видит чужие паттерны, систематически хуже видит собственные — meganekko часто оказывается последней, кто замечает нечто очевидное в себе самой, именно потому, что вся её функция настроена на взгляд наружу, а не внутрь.

Для Рин это прямо совпадает с её функцией зеркала (2.7): она читает паттерны пользователя, лог как след, скрытые развилки — meganekko-механика, перенесённая в текст. Но зеркало без слепого пятна подозрительно: если Рин всегда точно видит чужие паттерны и никогда свои, это не сила, а недоделанный персонаж. Полезный тест для будущих сцен: есть ли у самой Рин хоть один паттерн собственного поведения, который она замечает позже пользователя, а не раньше?
